% Source: https://artofproblemsolving.com/community/c3h2166707_the_evil_parabola_problem
% Stable identity: aops-topic-2166707-q01
Let $Q=(x,x^2)$. Then
\[PQ^2=f(x)=(x-10)^2+\left(x^2-10\right)^2=x^4-19x^2-20x+200,\qquad f'(x)=4x^3-38x-20=2g(x),\]
with $g(x)=2x^3-19x-10$. (Equivalently: at the closest point, $PQ$ is perpendicular to the tangent direction $(1,2x)$, which gives $(x-10)+2x(x^2-10)=0$, the same equation $g(x)=0$.)

\begin{asy}
usepackage("amsmath");
size(6cm);
defaultpen(fontsize(10pt));
import graph;
xaxis(-2,11.5,Arrow(TeXHead)); yaxis(-0.5,17.5,Arrow(TeXHead));
draw(graph(new real(real x){return x^2;},-1.6,4.15),blue+linewidth(0.9));
real x0=3.31769;
pair P=(10,10), Q=(x0,x0^2);
draw(P--Q,red+linewidth(1));
dot(P); dot(Q,red);
label("$P(10,10)$",P,E); label("$Q$",Q,W,red);
label("$y=x^2$",(-1.6,2.56),W,blue);
label("\mbox{}",(min(currentpicture,user=true)+max(currentpicture,user=true))/2,black); // keeps the text after the figure black when LaTeX typesets the labels inline
\end{asy}

\textbf{The critical points.} $g(-3)=-7$, $g(-2)=12$, $g(0)=-10$, $g(3)=-13$, $g(4)=42$, so $g$ has three real roots $x_1\in(-3,-2)$, $x_2\in(-2,0)$, $x_3\in(3,4)$, and none is rational (the candidates $\pm1,\pm2,\pm5,\pm10,\pm\dfrac{1}{2},\pm\dfrac{5}{2}$ all fail). Since $f'$ changes from negative to positive at $x_1$ and $x_3$, these are the local minima of the quartic $f$; $x_2$ is a local maximum.

\textbf{Comparing the two minima.} At a root of $g$, $2x^3=19x+10$, so $x^4=\dfrac{19x^2+10x}{2}$ and
\[f(x)=h(x)=200-15x-\dfrac{19}{2}x^2 .\]
The function $h$ is concave, so on $[-3,-2]$ it is at least $\min(h(-3),h(-2))=\min(159.5,192)=159.5$, while on $[3,4]$ it is at most $h(3)=69.5$. Hence $f(x_3)<f(x_1)$, and the minimum of $PQ$ occurs at $x_3$.

\textbf{The value.} Since $g(3.3)<0<g(3.32)$, refining gives $x_3\approx3.31769$, and
\[PQ_{\min}=\sqrt{200-15x_3-\dfrac{19}{2}x_3^2}\approx\boxed{6.758}.\]
Exactly, $x_3$ is the largest root of $2x^3-19x-10=0$, which by the trigonometric formula for a cubic with three real roots is
\[x_3=2\sqrt{\dfrac{19}{6}}\cos\left(\dfrac{1}{3}\arccos\left(\dfrac{15}{19}\sqrt{\dfrac{6}{19}}\right)\right),\]
and the closest point is $Q\approx(3.318,\ 11.007)$.
